\documentclass[french]{book}
\usepackage{teach}
\usepackage{verbatim}
\usepackage{amsmath,amsfonts,amssymb}
\usepackage{pgfplots}
\usepackage{mwe}
\usepackage{yhmath} 
%\usepackage{pstricks,pst-plot,pst-text,pst-tree,pst-eucl}
%
%\newcommand{\arc}[1]{\wideparen{#1}}
\RequirePackage{graphicx}
\RequirePackage{ifpdf}
 \ifpdf
   \DeclareGraphicsRule{*}{mps}{*}{}
 \fi
%
%\newcommand{\sysd}[2]{%
%       \left\{
%       \begin{aligned}
%          #1\\
%          #2\\
%       \end{aligned}
%       \right.%
%       }
%  \usepackage{fancyhdr}
%
%  \usepackage{amsmath}
%  \usepackage{amssymb}
%  \usepackage{amsfonts}
%  \usepackage{amsthm}
%  \usepackage{mathrsfs}
%
%  \usepackage{array}
  \usepackage{multicol}
%  \setlength{\multicolsep}{3pt}
%
  \usepackage{graphicx}
  \graphicspath{{Figures/}}
%  \usepackage{pstricks,pst-plot,pst-text,pst-tree,pst-eucl}
%  \usepackage[all]{xy}
%
%  \usepackage{fancybox}
%  \usepackage{color}
%
%%  \usepackage[frenchb]{babel}
%  \usepackage{hyperref}
%
%\usepackage{subfig}
%\pgfplotsset{compat=newest}
%\usepackage[all]{xy}
%\usepackage{geometry}
%\geometry{top=1cm, bottom=1cm, left=2cm, right=2cm}
%\usepackage[utf8]{inputenc}
%\usepackage[french]{babel}
\redefinecolor{thm}{title}{bgcolor}{alizarin}
\redefinecolor{prop}{title}{bgcolor}{meatbrown}
\redefinecolor{defn}{title}{bgcolor}{olivine}
\definecolor{alizarin}{rgb}{0.82, 0.1, 0.26}
\definecolor{meatbrown}{rgb}{0.9, 0.72, 0.23}
\definecolor{olivine}{rgb}{0.6, 0.73, 0.45}
%\usepackage{mathrsfs}
%%
%\newcommand{\R}{\mathbb {R}}
%\newcommand{\C}{\mathds {C}}
%\newcommand{\N}{\mathbb {N}}
%\newcommand{\Z}{\mathbb {Z}}
%\newcommand{\Q}{\mathbb {Q}}
%\newcommand{\D}{\mathbb {D}}
%%
%\newcommand{\se}{\geq}
%\newcommand{\ie}{\leq}
%\newcommand{\tm}{\times}
%\newcommand{\ell}{l}
%\newcommand{\ell'}{l'}
%%
%\newcommand{\ds}{\displaystyle}
%\newcommand{\limn}{\lim_{n\rightarrow +\infty}}
%\newcommand{\limo}[1][x]{\ds\lim_{#1\rightarrow 0}}
%\newcommand{\limite}[2][x]{\ds\lim_{#1\rightarrow #2}}
%\newcommand{\limpinf}[1][x]{\ds\lim_{#1\rightarrow +\infty}}
%\newcommand{\limminf}[1][x]{\ds\lim_{#1\rightarrow -\infty}}
%\newcommand{\limiteg}[2][x]{\ds\lim_{#1\xrightarrow{<}#2}}
%\newcommand{\limited}[2][x]{\ds\lim_{#1\xrightarrow{>}#2}}
%%
\newcommand{\vect}[1]{\overrightarrow{#1}}
\newcommand{\oij}{(\textit{O},{\overrightarrow{i}},{\overrightarrow{j}})}
%
\newcommand{\cp}[2]{%
        \begin{pmatrix}
        #1\\
        #2
        \end{pmatrix}%
        }
%
%\newcommand{\calb}{\mathscr{B}}
%\newcommand{\calc}{\mathscr{C}}
%\newcommand{\cald}{\mathscr{D}}
%\newcommand{\cale}{\mathscr{E}}
%\newcommand{\calf}{\mathscr{F}}
%\newcommand{\calp}{\mathscr{P}}
%\newcommand{\cals}{\mathscr{S}}
%\newcommand{\calt}{\mathscr{T}}
%
%\newcommand{\suite}[1][u]{\left( #1_{n} \right)_{n\in\N}}
%\newcommand{\suitep}[1][u]{\left( #1_{n} \right)_{n\in\N^{*}}}
%\usetikzlibrary{arrows}
%
%\newcommand{\poly}[1][a]{#1_0+#1_1x+\cdots+#1_nx^n}


\begin{document}
\chapter{Produit scalaire}
%\begin{minipage}{9cm}
Dans tout le chapitre, $\vect{u}$ et $\vect{v}$ d\'esignent deux vecteurs du plan.% et $O$, $A$ et $B$ sont trois points tels que $\vect{OA}=\vect{u}$ et $\vect{OB}=\vect{v}$.

Une unit\'e de longueur est fix\'ee.
% \end{minipage}
% \begin{minipage}{5cm}
% \begin{center}
% \includegraphics{ProdScal2.3}
% \end{center}
% \end{minipage}

\section{D\'efinition}
\begin{defn}
On appelle produit scalaire de $\vect{u}$ et $\vect{v}$ le nombre r\'eel not\'e $\vect{u}.\vect{v}$ et d\'efini par :
\[\vect{u}.\vect{v}=
\begin{cases}
 \Vert\vect{u}\Vert\times\Vert\vect{v}\Vert\times \cos\left(\vect{u};\vect{v}\right)&\text{si }\vect{u}\neq\vect{0}\text{ et }\vect{v}\neq\vect{0}\\
 0&\text{si }\vect{u}=\vect{0}\text{ ou }\vect{v}=\vect{0}
\end{cases}
\]
Le produit scalaire d'un vecteur $\vect{u}$ par lui-m\^eme ($\vect{u}.\vect{u}$) est appel\'e carr\'e scalaire de $\vect{u}$ et est \'egal à  $\Vert\vect{u}\Vert^2$ .
\end{defn}

\begin{rem}
Le produit scalaire est donc une op\'eration dont les arguments sont des vecteurs et dont le r\'esultat est r\'eel.
\end{rem}

\begin{minipage}{11cm}

\begin{exemple}
Sur la figure ci-contre, le triangle $OAB$ est \'equilat\'eral et $OA=2$. Calculer $\vect{u}.\vect{v}$.
\end{exemple}


On a $\Vert\vect{u}\Vert=2$,  $\Vert\vect{v}\Vert=2$ et $\left(\vect{u};\vect{v}\right)=\dfrac{\pi}{3}$

Ainsi $\vect{u}.\vect{v}=2\times2\times \cos\left( \dfrac{\pi}{3} \right)=2\times2\times \dfrac{1}{2}=2$

\end{minipage}
\hfill
\begin{minipage}{3.5cm}
\begin{center}
\includegraphics{CoursProdScalFigures.10}
\end{center}
\end{minipage}

\section{Autres expressions du produit scalaire}
\subsection{Cas des vecteurs colin\'eaires}
\begin{prop}
\begin{itemize}
\item Si $\vect{u}$ et $\vect{v}$ sont colin\'eaires et de m\^eme sens alors
\[\vect{u}.\vect{v}=\left\|\vect{u}\right\|\times \left\|\vect{v}\right\|\]
\item Si $\vect{u}$ et $\vect{v}$ sont colin\'eaires et de sens contraires alors \[\vect{u}.\vect{v}=-\Vert\vect{u}\Vert\times \Vert\vect{v}\Vert\]
\end{itemize}
\end{prop}

%\newpage
\begin{demo}
\begin{itemize}
\item Si $\vect{u}$ et $\vect{v}$ sont colin\'eaires et de m\^eme sens alors $\left(\vect{u};\vect{v}\right)=0$

Ainsi, $\vect{u}.\vect{v}=\left\|\vect{u}\right\|\times \left\|\vect{v}\right\|\times \cos0 = \left\|\vect{u}\right\| \times \left\|\vect{v}\right\|$

\item Si $\vect{u}$ et $\vect{v}$ sont colin\'eaires et de sens contraires, alors $\left(\vect{u};\vect{v}\right)=\pi$

Ainsi, $\vect{u}.\vect{v}=\left\|\vect{u}\right\|\times \left\|\vect{v}\right\|\times \cos\pi = -\left\|\vect{u}\right\| \times \left\|\vect{v}\right\|$
\end{itemize}
\end{demo}

Cons\'equences :
\begin{itemize}
 \item Si $\vect{u}$ et $\vect{v}$ sont colin\'eaires et de m\^eme sens alors $\vect{u}.\vect{v}>0$ et si $\vect{u}$ et $\vect{v}$ sont colin\'eaires et sens contraires alors $\vect{u}.\vect{v}<0$
 \item Quel que soit le vecteur $\vect{u}$, $\vect{u}^{2}=\Vert\vect{u}\Vert^{2}$
\end{itemize}

%
\subsection{Avec des projet\'es orthogonaux}
\begin{defn}
 On consid\`ere un vecteur $\vect{u}$ et une droite $d$. Si $A$ et $B$ sont deux points tels que $\vect{u}=\vect{AB}$ et $A'$ et $B'$ sont les projet\'es orthogonaux de $A$ et $B$ sur $d$ alors le vecteur $\vect{u'}=\vect{A'B'}$ ne d\'epend pas de $A$ et $B$ et est appel\'e projet\'e orthogonal de $\vect{u}$ sur $d$.
\end{defn}

\begin{prop}
Soient $A$, $B$, $C$ et $D$ quatre points du plan.

Si $C'$ et $D'$ sont les projet\'e orthogonaux de $C$ et $D$ sur $(AB)$ alors $\vect{AB}.\vect{CD}=\vect{AB}.\vect{C'D'}$
\end{prop}
\begin{center}
\includegraphics{CoursProdScalFigures.1}
\hfill
\includegraphics{CoursProdScalFigures.2}
\end{center}

\begin{demo}
Soit $E$ le point  tel que $\vect{C'E}=\vect{CD}$.

On a alors $\vect{AB}.\vect{CD}=\vect{AB}.\vect{C'E}=AB\times C'E\times \cos(\vect{AB};\vect{C'E})$.
\begin{itemize}
 \item Si l'angle $(\vect{AB};\vect{C'E})$ est droit alors $\cos(\vect{AB};\vect{C'E})=\cos(\vect{AB};\vect{CD})=0$ donc $\vect{AB}.\vect{CD}=\vect{AB}.\vect{C'D'}=0$.
 \item Si l'angle $(\vect{AB};\vect{C'E})$ est aigu alors $\vect{AB}$ et $\vect{C'D'}$ sont colin\'eaires et de m\^eme sens et $\cos(\vect{AB};\vect{C'E})=\dfrac{C'D'}{C'E}$
 
 Ainsi $\vect{AB}.\vect{CD}=AB\times C'E\times \dfrac{C'D'}{C'E}=AB\times C'D'=\vect{AB}.\vect{C'D'}$.
 \item Si l'angle $(\vect{AB};\vect{C'E})$ est obtus alors $\vect{AB}$ et $\vect{C'D'}$ sont colin\'eaires et de sens contraires et $\cos(\vect{AB};\vect{C'E})=-\dfrac{C'D'}{C'E}$
 
 Ainsi $\vect{AB}.\vect{CD}=-AB\times C'E\times \dfrac{C'D'}{C'E}=-AB\times C'D'=\vect{AB}.\vect{C'D'}$.
\end{itemize}
\end{demo}

\begin{minipage}{10cm}
\begin{exemple}
$ABC$ est un triangle isoc\`ele en $A$ tel que $AB=3$ et $BC=4$. $O$ est le milieu du segment $[BC]$. Calculer $\vect{BA}.\vect{BC}$ et $\vect{CA}.\vect{BC}$.
\end{exemple}

Le projet\'e orthogonal de $A$ sur $(BC)$ est $O$ donc

$\vect{BA}.\vect{BC}=\vect{BO}.\vect{BC}=BO\times BC=2\times3=6$

Le projet\'e orthogonal de $A$ sur $(BC)$ est $O$ donc

$\vect{CA}.\vect{BC}=\vect{CO}.\vect{BC}=-CO\times BC=-2\times3=-6$


\end{minipage}
\hfill
\begin{minipage}{4.5cm}
\begin{flushright}
\includegraphics{CoursProdScalFigures.8}
\end{flushright}
\end{minipage}

\newpage
\subsection{Dans un rep\`ere}
\begin{prop}
Soit $\oij$  un rep\`ere orthonormal du plan.

Si $\vect{u}\cp{x}{y}$ et $\vect{v}\cp{x'}{y'}$ alors $\vect{u}.\vect{v}=xx'+yy'$.
\end{prop}

\begin{demo}
Soit $A$ le point tel que $\vect{OA}=\vect{u}$ et $B$ le point tel que $\vect{OB}=\vect{v}$.
\begin{enumerate}
 \item Cas des vecteurs colin\'eaires.
 
 Il existe $k$ tel que $\vect{OB}=k\vect{OA}$. Ainsi $\vect{OB}\cp{kx}{ky}$.
 
   \begin{itemize}
    \item Si $\vect{OA}$ et $\vect{OB}$ sont colin\'eaires et de m\^eme sens alors $k>0$ et $OB=kOA$.

    On a alors $\vect{OA}.\vect{OB}=OA\times OB=k\times OA\times OA=kOA^2$.
    \item Si $\vect{OA}$ et $\vect{OB}$ sont colin\'eaires et de m\^eme sens contraire alors $k<0$ et $OB=-kOA$.

    On a alors $\vect{OA}.\vect{OB}=-OA\times OB=k\times OA\times OA=kOA^2$.
   \end{itemize}
  Ainsi $\vect{OA}.\vect{OB}=k(x^2+y^2)=kxx+kyy=xx'+yy'$
 \item Cas des vecteurs quelconques.
 
 \begin{minipage}{9cm}
 Soit $H$ le projet\'e orthogonal de $B$ sur $(OA)$.
 
 On a alors, d'apr\`es ce qui pr\'ec\`ede, \[\vect{OA}.\vect{OB}=\vect{OA}.\vect{OH}=xx_H+yy_H\]
 
 En appliquant le th\'eor\`eme de Pythagore dans les triangles $OBH$ et $ABH$ rectangles en $H$, on obtient :
 \[BH^2=OB^2-OH^2=AB^2-AH^2\]
 \end{minipage}
 \hfill
 \begin{minipage}{3.5cm}
  \begin{center}
   \includegraphics{CoursProdScalFigures.4}
  \end{center}
 \end{minipage}

 ce qui nous donne :
 \[x'^2+y'^2-x_H^2-y_H^2=(x-x')^2+(y-y')^2-(x-x_H)^2-(y-y_H)^2\]
 puis, en simplifiant :
 \[0=-2xx'-2yy'+2xx_H+2yy_H\]
 soit :
 \[xx_H+yy_H=xx'+yy'\]
 On en d\'eduit donc : $\vect{OA}.\vect{OB}=xx'+yy'$.
\end{enumerate}
\end{demo}

\begin{exemple}
Dans un rep\`ere orthonormal $\oij$, $\vect{u}\cp{6}{3}$, $\vect{v}\cp{3}{-1}$ et $\vect{w}\cp{-2}{2}$

Calculer $\vect{u}.\vect{v}$, $\vect{u}.\vect{w}$ et $\vect{v}.\vect{w}$.
\end{exemple}


$\vect{u}.\vect{v}= 6\times3+3\times(-1)=18-3=15$

$\vect{u}.\vect{w}=6\times(-2)+3\times2=-12+6=-6$

$\vect{v}.\vect{w}=3\times(-2)+(-1)\times2=-6-2=-8$



\section{R\`egles de calcul}
\begin{prop}
Quels que soient $\vect{u}$ et $\vect{v}$ :

\vspace{0.05cm}
\begin{itemize}
\begin{minipage}[c]{11cm}
\item $\vect{0}.\vect{v}=\vect{u}.\vect{0}=0$
\vspace{0.2cm}
\item $\vect{u}.\vect{v}=\vect{v}.\vect{u}$ \hfill (On dit que le produit scalaire est sym\'etrique)
\vspace{0.2cm}
\end{minipage}

\begin{minipage}{9cm}
\item $\vect{u}.(\vect{v}+\vect{w})=\vect{u}.\vect{v}+\vect{u}.\vect{w}$ et $\vect{u}.(k\vect{v})=k\times (\vect{u}.\vect{v})$

\vspace{0.2cm}
\item $(\vect{u}+\vect{v}).\vect{w}=\vect{u}.\vect{w}+\vect{v}.\vect{w}$ et $(k\vect{u}).\vect{v}=k\times (\vect{u}.\vect{v})$
\end{minipage}
\hspace{0.5cm}
\begin{minipage}[c]{4cm}
(On dit que le produit scalaire est bilin\'eaire)
\end{minipage}
\end{itemize}
\end{prop}

\begin{demo}
\begin{itemize}
\item[$\bullet$] Le premier r\'esultat est une cons\'equence directe de la d\'efinition.
\item[$\bullet$] Le deuxi\`eme r\'esultat est aussi une cons\'equence de la d\'efinition ($\cos(\vect{u};\vect{v})=\cos(\vect{v};\vect{u})$).
\item[$\bullet$] On se place dans un rep\`ere orthonormal du plan et on pose $\vect{u}\cp{x}{y}$, $\vect{v}\cp{x'}{y'}$ et $\vect{w}\cp{x''}{y''}$

Les coordonn\'ees de $\vect{v}+\vect{w}$ sont $\cp{x'+x''}{y'+y''}$.

On a ainsi, $\vect{u}.(\vect{v}+\vect{w})=x(x'+x'')+y(y'+y'')=xx'+xx''+yy'+yy''$.

De plus $\vect{u}.\vect{v}+\vect{u}.\vect{w}=xx'+yy'+xx''+yy''$.

Conclusion : $\vect{u}.(\vect{v}+\vect{w})=\vect{u}.\vect{v}+\vect{u}.\vect{w}$

Les coordonn\'ees de $k\vect{v}$ sont $\cp{kx'}{ky'}$.

On a ainsi $\vect{u}.(k\vect{v}) = x(kx')  + y(ky') = kxx'  + kyy'$.

De plus $k\times (\vect{u}.\vect{v})=k(xx'+yy')=kxx'+kyy'$

Conclusion $\vect{u}.(k\vect{v})=k\times (\vect{u}.\vect{v})$

\item[$\bullet$] Le quatri\`eme r\'esultat se d\'emontre en utilisant les deux pr\'ec\'edents.
\end{itemize}
\end{demo}

\newpage
\begin{exemple}
$\vect{u}$ et $\vect{v}$ sont deux vecteurs, simplifier $\left(\vect{u}+\vect{v}\right)^{2}$, $\left(\vect{u}-\vect{v}\right)^{2}$ et $\left(\vect{u}+\vect{v}\right).\left(\vect{u}-\vect{v}\right)$.

\begin{multicols}{2}

\begin{align*}
\begin{split}
\left(\vect{u}+\vect{v}\right)^{2}&=\left(\vect{u}+\vect{v}\right).\left(\vect{u}+\vect{v}\right)\\
&=\left(\vect{u}+\vect{v}\right).\vect{u}+\left(\vect{u}+\vect{v}\right).\vect{v}\\
&=\vect{u}^2 + \vect{v}.\vect{u} + \vect{v}.\vect{u} + \vect{v}^2 \\
&= \vect{u}^2 +  2\vect{u}.\vect{v} + \vect{v}^2
\end{split}
\end{align*}

\begin{align*}
\begin{split}
\left(\vect{u}-\vect{v}\right)^{2}&=\left(\vect{u}-\vect{v}\right).\left(\vect{u}-\vect{v}\right)\\
&=\left(\vect{u}-\vect{v}\right).\vect{u}-\left(\vect{u}-\vect{v}\right).\vect{v}\\
&=\vect{u}^2 - \vect{v}.\vect{u} - \vect{v}.\vect{u} + \vect{v}^2 \\
&= \vect{u}^2 -  2\vect{u}.\vect{v} + \vect{v}^2
\end{split}
\end{align*}
\end{multicols}

\begin{align*}
\begin{split}
\left(\vect{u}+\vect{v}\right).\left(\vect{u}-\vect{v}\right)&=\vect{u}^2 - \vect{u}.\vect{v} + \vect{v}.\vect{u} - \vect{v}^2 \\&= \vect{u}^2 - \vect{v}^2
\end{split}
\end{align*}
\end{exemple}

\section{Vecteurs orthogonaux}
\subsection{D\'efinition}
\begin{defn}
Les vecteurs $\vect{u}$ et $\vect{v}$ sont orthogonaux si :
\[\vect{u}=\vect{0} \; \text{ou} \; \vect{v}=\vect{0} \; \text{ou} \; (OA)\bot(OB)\]
\end{defn}

\subsection{Propri\'et\'e}
\begin{prop}
Les vecteurs $\vect{u}$ et $\vect{v}$ sont orthogonaux si et seulement si $\vect{u}.\vect{v}=0$.
\end{prop}

\begin{demo}
\vspace{-1em}
\begin{align*}
 \vect{u}.\vect{v}=0 & \Leftrightarrow \left\|\vect{u}\right\| \times \left\|\vect{v}\right\| \times \cos(\vect{u};\vect{v})=0\\
 & \Leftrightarrow \left\|\vect{u}\right\|=0 \text{ ou } \left\|\vect{v}\right\|=0 \text{ ou } \cos(\vect{u};\vect{v})=0\\
 & \Leftrightarrow \vect{u}=\vect{0} \text{ ou } \vect{v}=\vect{0} \text{ ou } (\vect{u};\vect{v})=\dfrac{\pi}{2}+k\times 2\pi\\
 & \Leftrightarrow \vect{u}\text{ et } \vect{v} \text{ sont orthogonaux}
\end{align*}

%\vspace{+1em}
\end{demo}

\begin{exemple}
Dans un rep\`ere orthonormal $\oij$, on donne $\vect{u}\cp{6}{4}$, $\vect{v}\cp{6}{-9}$.

D\'emontrer que $\vect{u}$ et $\vect{v}$ sont orthogonaux.
\end{exemple}


$\vect{u}.\vect{v}= 6\times6  + 4\times(-9)  = 36 - 36 =  0$.
Les vecteurs $\vect{u}$ et $\vect{v}$ sont donc orthogonaux.





%%%%%%%%%%%         Ancienne version

%%
%%
% \begin{minipage}{9cm}
% Dans tout le chapitre, $\vect{u}$ et $\vect{v}$ d\'esignent deux vecteurs du plan et $O$, $A$ et $B$ sont trois points tels que $\vect{OA}=\vect{u}$ et $\vect{OB}=\vect{v}$.
% 
% Une unit\'e de longueur est fix\'ee.
%  \end{minipage}
%  \begin{minipage}{5cm}
%  \begin{center}
%  \includegraphics{ProdScal2.3}
%  \end{center}
%  \end{minipage}
%  
% \section{D\'efinition}
% \label{defps}
% \begin{definition}
% On appelle produit scalaire de $\vect{u}$ et $\vect{v}$ le nombre r\'eel not\'e $\vect{u}.\vect{v}$ et d\'efini par :
% \[\vect{u}.\vect{v}=\dfrac{1}{2}\left(\Vert\vect{u}\Vert^{2}+\Vert\vect{v}\Vert^{2}-\Vert\vect{v}-\vect{u}\Vert^{2}\right)\]
% Le produit scalaire d'un vecteur $\vect{u}$ par lui-m\^eme ($\vect{u}.\vect{u}$) est appel\'e carr\'e scalaire de $\vect{u}$ et se note $\vect{u}^{2}$. On a donc $\vect{u}^{2}=\Vert\vect{u}\Vert^{2}$
% \end{definition}
% Remarques :
% \begin{itemize}
% \item Le produit scalaire est donc une op\'eration dont les arguments sont des vecteurs et dont le r\'esultat est r\'eel.
% \item Si un des vecteurs est nul alors le produit scalaire est nul.
% \end{itemize}
% 
% 
% \begin{minipage}{11cm}
% \begin{ex}
% Sur la figure ci-contre, le triangle $OAB$ est \'equilat\'eral et $OA=2$. Calculer $\vect{u}.\vect{v}$.
% \end{ex}
% 
% 
% On a $\Vert\vect{u}\Vert=2$,  $\Vert\vect{v}\Vert=2$ et $\Vert\vect{v}-\vect{u}\Vert=\Vert\vect{OB}-\vect{OA}\Vert=\Vert\vect{AB}\Vert=2$
% 
% Ainsi $\vect{u}.\vect{v}=\dfrac{1}{2}\times(2^2+2^2-2^2)=2$
% 
% \end{minipage}
% \hfill
% \begin{minipage}{3.5cm}
% \begin{center}
% \includegraphics{ProdScal2.6}
% \end{center}
% \end{minipage}
% %%
% %%
% \section{Autres expressions du produit scalaire}
% \subsection{Cas des vecteurs colin\'eaires}
% \begin{prop}
% \begin{itemize}
% \item Si $\vect{u}$ et $\vect{v}$ sont colin\'eaires et de m\^eme sens alors
% \[\vect{u}.\vect{v}=\left\|\vect{u}\right\|\times \left\|\vect{v}\right\|=OA\times OB\]
% \item Si $\vect{u}$ et $\vect{v}$ sont colin\'eaires et de sens contraires alors \[\vect{u}.\vect{v}=-\Vert\vect{u}\Vert\times \Vert\vect{v}\Vert=-OA\times OB\]
% \end{itemize}
% \end{prop}
% 
% %\newpage
% \begin{demo}
% Lorsque $\vect{u}$ et $\vect{v}$ sont colin\'eaires, on peut exprimer $\left\|\vect{v}-\vect{u}\right\|$ en fonction de $\left\|\vect{v}\right\|$ et $\left\|\vect{u}\right\|$.
% \begin{itemize}
% \item Si $\vect{u}$ et $\vect{v}$ sont colin\'eaires et de m\^eme sens alors $\left\|\vect{v}-\vect{u}\right\|^2=\left( \left\|\vect{v}\right\|-\left\|\vect{u}\right\| \right)^2$.
% 
% On a ainsi,
% %\begin{align}
% \[\vect{u}.\vect{v}=\dfrac{1}{2}\left(\left\|\vect{u}\right\|^{2}+\left\|\vect{v}\right\|^{2}-\left\|\vect{u}\right\|^2+2\left\|\vect{u}\right\|\left\|\vect{v}\right\|-\left\|\vect{u}\right\|^{2}\right)=\left\|\vect{u}\right\|\times \left\|\vect{v}\right\|\]
% %\end{align}
% \item Si $\vect{u}$ et $\vect{v}$ sont colin\'eaires et de sens contraires, % alors
%  $\left\|\vect{v}-\vect{u}\right\|^2=\left( \left\|\vect{v}\right\|+\left\|\vect{u}\right\| \right)^2$.
% 
% On a ainsi,
% %\begin{align}
% \[\vect{u}.\vect{v}=\dfrac{1}{2}\left(\left\|\vect{u}\right\|^{2}+\left\|\vect{v}\right\|^{2}-\left\|\vect{u}\right\|^2-2\left\|\vect{u}\right\|\left\|\vect{v}\right\|-\left\|\vect{u}\right\|^{2}\right)=-\left\|\vect{u}\right\|\times \left\|\vect{v}\right\|\]
% \end{itemize}
% \end{demo}
% 
% Cons\'equences : Si $\vect{u}$ et $\vect{v}$ sont colin\'eaires et de m\^eme sens alors $\vect{u}.\vect{v}>0$ et si $\vect{u}$ et $\vect{v}$ sont colin\'eaires et sens contraires alors $\vect{u}.\vect{v}<0$
% %
% \subsection{Avec un projet\'e orthogonal}
% \begin{prop}
% Si $H$ est le projet\'e orthogonal de $B$ sur $(OA)$ alors $\vect{u}.\vect{v}=\vect{OA}.\vect{OH}$.
% \end{prop}
% \begin{center}
% \includegraphics{ProdScal2.1}
% \hspace{2cm}
% \includegraphics{ProdScal2.2}
% \end{center}
% 
% \begin{dem}
% Les triangles $OBH$ et $ABH$ sont rectangles en $H$. On a donc, d'apr\`es le th\'eor\`eme de Pythagore, $OB^2=OH^2+HB^2$ et $AB^2=AH^2+HB^2$.
% Ainsi :
% \begin{align*}
% \vect{u}.\vect{v} &= \dfrac{1}{2}\left( OA^2+OB^2-AB^2 \right) = \dfrac{1}{2}\left( OA^2+OH^2+HB^2-AH^2-HB^2 \right)\\
%  &= \dfrac{1}{2}\left(OA^2+OH^2-AH^2\right) = \dfrac{1}{2}\left(\left\|\vect{OA}\right\|^2 +  \left\|\vect{OB}\right\|^2 - \left\|\vect{OH}-\vect{OA}\right\|^2\right)\\
% &= \vect{OA}.\vect{OH}
% \end{align*}
% 
% \end{dem}
% 
% \begin{minipage}{10cm}
% \begin{ex}
% $ABC$ est un triangle isoc\`ele en $A$ tel que $AB=3$ et $BC=4$. $O$ est le milieu du segment $[BC]$. Calculer $\vect{BA}.\vect{BC}$ et $\vect{CA}.\vect{BC}$.
% \end{ex}
% 
% Le projet\'e orthogonal de $A$ sur $(BC)$ est $O$ donc
% 
% $\vect{BA}.\vect{BC}=\vect{BO}.\vect{BC}=BO\times BC=2\times3=6$
% 
% Le projet\'e orthogonal de $A$ sur $(BC)$ est $O$ donc
% 
% $\vect{CA}.\vect{BC}=\vect{CO}.\vect{BC}=-CO\times BC=-2\times3=-6$
% 
% 
% \end{minipage}
% \hfill
% \begin{minipage}{4.5cm}
% \begin{flushright}
% \includegraphics{ProdScal2.4}
% \end{flushright}
% \end{minipage}
% %
% \subsection{Avec la fonction cosinus}
% \begin{prop}
% Quels que soient les points $O$, $A$ et $B$,
% \[\vect{OA}.\vect{OB}=OA\times OB\times \cos\left(\widehat{AOB}\right)\]
% \end{prop}
% 
% \begin{dem}
% Soit $H$ le projet\'e orthogonal de $B$ sur $(OA)$. D'apr\`es la propri\'et\'e pr\'ec\'edente, $\vect{OA}.\vect{OB}=\vect{OA}.\vect{OH}$.
% \begin{itemize}
% \item Si $\widehat{AOB}$ est un angle droit alors $\vect{OH}=\vect{0}$ donc $\vect{OA}.\vect{OB}=0$. De plus le cosinus d'un angle droit \'etant nul, on a bien $\vect{OA}.\vect{OB}=OA\times OB\times \cos(\widehat{AOB})$.
% \item Si $\widehat{AOB}$ est un angle aigu alors $OA$ et $OH$ sont colin\'eaires et de m\^eme sens donc $\vect{OA}.\vect{OB}=OA\times OH$.
% 
% De plus, $OH=OB\cos( \widehat{HOB} )=OB\cos( \widehat{AOB} )$
% 
% donc $\vect{OA}.\vect{OB}=OA\times OB\times \cos(\widehat{AOB})$
% 
% \item Si $\widehat{AOB}$ est un angle obtus alors $OA$ et $OH$ sont colin\'eaires et de sens contraires donc $\vect{OA}.\vect{OB}=-OA\times OH$.
% 
% De plus, $OH=OB\cos( \widehat{HOB})$ et $\widehat{HOB}$ et $\widehat{AOB}$ sont suppl\'ementaires donc $\cos( \widehat{HOB})=-\cos( \widehat{AOB})$
% 
% Ainsi $\vect{OA}.\vect{OB}=-OA\times OB\times -\cos(\widehat{AOB})=OA\times OB\times \cos(\widehat{AOB})$
% \end{itemize}
% 
% \end{dem}
% 
% \begin{minipage}{10cm}
% \begin{ex}
% $ABC$ est un triangle \'equilat\'eral inscrit dans un cercle $\calc$ de centre $O$ et de rayon 3. Calculer $\vect{OB}.\vect{OC}$.
% \end{ex}
% 
% 
% $\vect{OB}.\vect{OC}=OB\times OC\times \cos\left( \vect{OB},\vect{OC}\right)=3\times3\times\cos(120\text{\degres})=-\dfrac{9}{2}$
% 
% \end{minipage}
% \hfill
% \begin{minipage}{4cm}
% \begin{flushright}
% \includegraphics{ProdScal2.5}
% \end{flushright}
% \end{minipage}
% %
% \subsection{Dans un rep\`ere}
% \begin{prop}
% Soit $\oij$  un rep\`ere orthonormal du plan.
% 
% Si $\vect{u}\cp{x}{y}$ et $\vect{v}\cp{x'}{y'}$ alors $\vect{u}.\vect{v}=xx'+yy'$.
% \end{prop}
% 
% \begin{dem}
% Si $\oij$ est un rep\`ere orthonormal alors $\left\|\vect{u}\right\|^2=x^2+y^2$, $\left\|\vect{v}\right\|^2=x'^2+y'^2$ et $\left\|\vect{v}-\vect{u}\right\|^2=(x'-x)^2+(y'-y)^2$.
% 
% On a alors :
% \begin{align*}
% \vect{u}.\vect{v}&=\dfrac{1}{2}(x^2+y^2+x'^2+y'^2-((x'-x)^2+(y'-y)^2))\\
% &=\dfrac{1}{2}(x^2+y^2+x'^2+y'^2-x'^2+2xx'-x^2-y'^2+2yy'-y^2)\\
% &=\dfrac{1}{2}(2xx'+2yy')=xx'+yy'
% \end{align*}
% \end{dem}
% 
% \begin{ex}
% Dans un rep\`ere orthonormal $\oij$, $\vect{u}\cp{6}{3}$, $\vect{v}\cp{3}{-1}$ et $\vect{w}\cp{-2}{2}$
% 
% Calculer $\vect{u}.\vect{v}$, $\vect{u}.\vect{w}$ et $\vect{v}.\vect{w}$.
% \end{ex}
% 
% 
% $\vect{u}.\vect{v}= 6\times3+3\times(-1)=18-3=15$
% 
% $\vect{u}.\vect{w}=6\times(-2)+3\times2=-12+6=-6$
% 
% $\vect{v}.\vect{w}=3\times(-2)+(-1)\times2=-6-2=-8$
% 
% %%
% %%
% \section{R\`egles de calcul}
% \begin{prop}
% Quels que soient $\vect{u}$ et $\vect{v}$ :
% 
% \begin{itemize}
% \begin{minipage}[c]{11cm}
% \item $\vect{0}.\vect{v}=\vect{u}.\vect{0}=0$
% \vspace{0.2cm}
% \item $\vect{u}.\vect{v}=\vect{v}.\vect{u}$ (On dit que le produit scalaire est sym\'etrique)
% \vspace{0.2cm}
% \end{minipage}
% 
% \begin{minipage}{9cm}
% \item $\vect{u}.(\vect{v}+\vect{w})=\vect{u}.\vect{v}+\vect{u}.\vect{w}$ et $\vect{u}.(k\vect{v})=k\times (\vect{u}.\vect{v})$
% 
% %\vspace{-0.3cm}
% \item $(\vect{u}+\vect{v}).\vect{w}=\vect{u}.\vect{w}+\vect{v}.\vect{w}$ et $(k\vect{u}).\vect{v}=k\times (\vect{u}.\vect{v})$
% \end{minipage}
% \hspace{0.5cm}
% \begin{minipage}[c]{4cm}
% (On dit que le produit scalaire est bilin\'eaire)
% \end{minipage}
% \end{itemize}
% \end{prop}
% 
% \begin{dem}
% \begin{itemize}
% \item Le premier r\'esultat est une cons\'equence directe de la d\'efinition (vu au paragraphe \ref{defps}).
% \item Le deuxi\`eme r\'esultat est aussi une cons\'equence de la d\'efinition ($\left\|\vect{v}-\vect{u}\right\|=\left\|\vect{u}-\vect{v}\right\|$).
% \item On se place dans un rep\`ere orthonormal du plan et on pose $\vect{u}\cp{x}{y}$, $\vect{v}\cp{x'}{y'}$ et $\vect{w}\cp{x''}{y''}$
% 
% Les coordonn\'ees de $\vect{v}+\vect{w}$ sont $\cp{x'+x''}{y'+y''}$.
% 
% On a ainsi, $\vect{u}.(\vect{v}+\vect{w})=x(x'+x'')+y(y'+y'')=xx'+xx''+yy'+yy''$.
% 
% De plus $\vect{u}.\vect{v}+\vect{u}.\vect{w}=xx'+yy'+xx''+yy''$.
% 
% Conclusion : $\vect{u}.(\vect{v}+\vect{w})=\vect{u}.\vect{v}+\vect{u}.\vect{w}$
% 
% Les coordonn\'ees de $k\vect{v}$ sont $\cp{kx'}{ky'}$.
% 
% On a ainsi $\vect{u}.(k\vect{v}) = x(kx')  + y(ky') = kxx'  + kyy'$.
% 
% De plus $k\times (\vect{u}.\vect{v})=k(xx'+yy')=kxx'+kyy'$
% 
% Conclusion $\vect{u}.(k\vect{v})=k\times (\vect{u}.\vect{v})$
% 
% \item Le quatri\`eme r\'esultat se d\'emontre en utilisant les deux pr\'ec\'edents.
% 
% \end{itemize}
% \end{dem}
% 
% \begin{ex}
% $\vect{u}$ et $\vect{v}$ sont deux vecteurs, simplifier $\left(\vect{u}+\vect{v}\right)^{2}$, $\left(\vect{u}-\vect{v}\right)^{2}$ et $\left(\vect{u}+\vect{v}\right).\left(\vect{u}-\vect{v}\right)$.
% \end{ex}
% 
% 
% \begin{align*}
% \begin{split}
% \left(\vect{u}+\vect{v}\right)^{2}&=\left(\vect{u}+\vect{v}\right).\left(\vect{u}+\vect{v}\right)=\left(\vect{u}+\vect{v}\right).\vect{u}+\left(\vect{u}+\vect{v}\right).\vect{v}\\
% &=\vect{u}^2 + \vect{v}.\vect{u} + \vect{v}.\vect{u} + \vect{v}^2 = \vect{u}^2 +  2\vect{u}.\vect{v} + \vect{v}^2
% \end{split}\\
% %\end{align*}
% %\begin{multline*}
% \begin{split}
% \left(\vect{u}-\vect{v}\right)^{2}&=\left(\vect{u}-\vect{v}\right).\left(\vect{u}-\vect{v}\right)=\left(\vect{u}-\vect{v}\right).\vect{u}-\left(\vect{u}-\vect{v}\right).\vect{v}\\
% &=\vect{u}^2 - \vect{v}.\vect{u} - \vect{v}.\vect{u} + \vect{v}^2 = \vect{u}^2 -  2\vect{u}.\vect{v} + \vect{v}^2
% \end{split}\\
% %\end{multline*}
% \left(\vect{u}+\vect{v}\right).\left(\vect{u}-\vect{v}\right)&=\vect{u}^2 - \vect{u}.\vect{v} + \vect{v}.\vect{u} - \vect{v}^2 = \vect{u}^2 - \vect{v}^2
% \end{align*}
% 
% 
% %%
% %%
% \section{Vecteurs orthogonaux}
% \subsection{D\'efinition}
% \begin{definition}
% Les vecteurs $\vect{u}$ et $\vect{v}$ sont orthogonaux si :
% \[\vect{u}=\vect{0} \; \text{ou} \; \vect{v}=\vect{0} \; \text{ou} \; (OA)\bot(OB)\]
% \end{definition}
% 
% \subsection{Propri\'et\'e}
% \begin{prop}
% Les vecteurs $\vect{u}$ et $\vect{v}$ sont orthogonaux si et seulement si $\vect{u}.\vect{v}=0$.
% \end{prop}
% 
% \begin{dem}
% Avec $\vect{OA}=\vect{u}$ et $\vect{OB}=\vect{v}$, on a : $\vect{u}.\vect{v}=\dfrac{1}{2}(OA^2 + OB^2 - AB^2)$.
% 
% \begin{itemize}
% \item Supposons que $\vect{u}$ et $\vect{v}$ sont orthogonaux.
% 
% Soit un des deux vecteurs est nul auquel cas $\vect{u}.\vect{v}=0$
% 
% Soit le triangle $OAB$ et rectangle en $O$  et  donc d'apr\`es le th\'eor\`eme de Pythagore, $OA^2 + OB^2 = AB^2$  et ainsi, $\vect{u}.\vect{v}=0$.
% 
% \item Supposons $\vect{u}.\vect{v}=0$.
% 
% Soit un des deux vecteurs est nul et donc $\vect{u}$ et $\vect{v}$ sont orthogonaux.
% 
% Soit $OA^2 + OB^2 = AB^2$ et d'apr\`es la r\'eciproque du th\'eor\`eme de Pythagore, $OAB$ est rectangle en $O$ donc $\vect{u}$ et $\vect{v}$ sont orthogonaux.
% \end{itemize}
% 
% \end{dem}
% 
% \begin{ex}
% Dans un rep\`ere orthonormal $\oij$, on donne $\vect{u}\cp{6}{4}$, $\vect{v}\cp{6}{-9}$.
% 
% D\'emontrer que $\vect{u}$ et $\vect{v}$ sont orthogonaux.
% \end{ex}
% 
% 
% $\vect{u}.\vect{v}= 6\times6  + 4\times(-9)  = 36 - 36 =  0$.
% Les vecteurs $\vect{u}$ et $\vect{v}$ sont donc orthogonaux.
% 
% 

\end{document}